\section{Totally Inseparable Extensions}
For degree checks involving two independent indeterminates, one useful elementary rule is
\[
[\mathbb F_3(y,z):\mathbb F_3(y^az^b,y^cz^d)]=|ad-bc|
\]
when the two exponent vectors are independent. Here is why the determinant appears. The subgroup of $\Z^2$ generated by $(a,b),(c,d)$ has $|ad-bc|$ cosets. Choose one monomial for each exponent coset. Every Laurent monomial is a coefficient-field monomial times one chosen representative, so these representatives span; their span is a finite-dimensional domain over the coefficient field, hence is a field and includes rational functions too. They are independent because clearing coefficient denominators leaves terms in distinct exponent cosets, which cannot cancel. Thus they form a basis. This supplies the degree certificates below; separability and the power tests establish the actual closure fields.

\exercise{52.1}{A separable closure with two independent exponents}
Let $K=\Z_3(y^{12},z^{18})$ and $L=\Z_3(y,z)$. Set $S=\Z_3(y^3,z^9)$. The generators $y^3$ and $z^9$ satisfy respectively $X^4-y^{12}$ and $X^2-z^{18}$ over $K$, whose nonzero roots are simple in characteristic 3. Thus $S/K$ is separable. Also $y^3,z^9\in S$, so every element of $L$ has some $3$-power in $S$; in fact its ninth power is in $S$. Hence $L/S$ is purely inseparable. An element separable over $K$ is separable over $S$ too, since its new minimal polynomial divides its old separable one; if it also lies in the purely inseparable extension $L/S$, it must be in $S$. Therefore $S$ is exactly the separable closure. The monomial degree counts are $[L:K]=12\cdot18=216$, $[L:S]=3\cdot9=27$, and $[S:K]=8$.

\exercise{52.2}{A separable closure with a mixed monomial}
Now $K=\Z_3(y^{12},y^2z^{18})$. Put $S=\Z_3(y^3,yz^9)$. The first generator satisfies the separable equation $X^4-y^{12}=0$ and the second the separable equation $X^2-y^2z^{18}=0$, so $S/K$ is separable. Moreover $y^3\in S$ and $z^{27}=(yz^9)^3/y^3\in S$. Thus $L=\Z_3(y,z)$ is purely inseparable over $S$. The separable/purely-inseparable intersection argument from Exercise 52.1 makes $S$ the full separable closure. For a degree check, the exponent lattice of $K$ has determinant $12\cdot18=216$, while that of $S$, generated by $(3,0),(1,9)$, has determinant 27; hence $[S:K]=8$.

\exercise{52.3}{A totally inseparable closure in the independent case}
The field is $P=\Z_3(y^4,z^2)$. Its generators have a third and a ninth power, respectively, in $K=\Z_3(y^{12},z^{18})$, so $P/K$ is purely inseparable. The extension $L/P$ is separable, because $y$ satisfies $X^4-y^4$ and $z$ satisfies $X^2-z^2$, both separable polynomials. If $a\in L$ has a $3$-power in $K$, it is purely inseparable over $P$ but also separable over $P$, so $a\in P$. This proves maximality. Thus the totally inseparable closure is $P$; its degree over $K$ is 27.

\exercise{52.4}{A totally inseparable closure in the mixed case}
Take $P=\Z_3(y^4,y^2z^2)$. It contains $K=\Z_3(y^{12},y^2z^{18})$, since $y^{12}=(y^4)^3$ and $y^2z^{18}=(y^2z^2)^9/(y^4)^4$. Its generators have $3$-powers in $K$:
\[(y^4)^3=y^{12},\qquad(y^2z^2)^{27}=(y^{12})^4(y^2z^{18})^3.\]
Hence $P/K$ is purely inseparable. Over $P$, first adjoining $y$ is separable via its fourth-power equation, and then adjoining $z$ is separable via $z^2=(y^2z^2)/y^2$. Thus $L/P$ is separable. The same intersection argument as in Exercise 52.3 shows that every element of $L$ purely inseparable over $K$ lies in $P$. Its exponent lattice, generated by $(4,0),(2,2)$, has determinant 8, so $[P:K]=216/8=27$.

\exercise{52.5}{Ten statements about totally inseparable extensions}
In this section Fraleigh reserves ``totally inseparable'' for a \emph{proper} extension with one embedding. The more common ``purely inseparable'' convention also permits the trivial extension.
\begin{enumerate}[label=(\alph*)]
\item False. $\Z_3(t)(\sqrt t)/\Z_3(t)$ is a proper separable extension, since $X^2-t$ is irreducible and has nonzero derivative.
\item True. A totally inseparable algebraic element has some $p^t$-power in $F$.
\item False. $\Z_5(y)/\Z_5(y^5)$ has degree 5 and only one embedding; its minimal polynomial is $X^5-y^5$.
\item False. $\Z_5(y)/\Z_5(y^{10})$ has inseparable degree 5 and separable degree 2, since $X^{10}-y^{10}$ has two distinct roots, each repeated five times.
\item False. The same mixed extension has two embeddings, so is not totally inseparable either.
\item False. The element $y$ over $\Z_5(y^{10})$ in (d) is neither separable nor totally inseparable.
\item True. The separable elements form a subfield by Exercise 51.13.
\item False with the book's strict convention: if $E/F$ is separable, its separable closure is $E$, and $E/E$ is not a proper extension. The corrected claim allows $E$ to equal that closure, or uses the inclusive ``purely inseparable'' convention.
\item True. If $E/F$ is not separable, its separable closure is proper, and every element outside it is totally inseparable over it.
\item True, by the minimal-polynomial characterization of total inseparability.
\end{enumerate}

\exercise{52.6}{Purely inseparable elements form a field}
Include the elements of $F$ as requested. For $a,b$ in the proposed set, choose a common power $q=p^m$ with $a^q,b^q\in F$. Then $(a+b)^q=a^q+b^q\in F$, $(ab)^q=a^qb^q\in F$, and if $b\ne0$, $(a/b)^q=a^q/b^q\in F$. Additive inverses work in the same way. Thus it is closed under field operations and contains 0,1. Every element outside $F$ with such a power is totally inseparable, by the power criterion, so these closure calculations prove the assertion. In characteristic zero no proper totally inseparable elements exist, and the set is simply $F$.

\exercise{52.7}{Perfection and surjective Frobenius}
If $F$ is perfect and $a\in F$, choose a root $b$ of $X^p-a$ in a closure. That polynomial equals $(X-b)^p$, so $b$ has only one conjugate. Separability forces its degree to be 1; hence $b\in F$ and $a=b^p$. Thus $F^p=F$.
Conversely suppose $F^p=F$. An inseparable irreducible polynomial would have derivative zero, so have form $\sum a_jX^{pj}$. Write $a_j=b_j^p$. Then it equals $(\sum b_jX^j)^p$, contradicting irreducibility. Therefore all irreducibles are separable, and $F$ is perfect.

\exercise{52.8}{Perfection through a finite extension}
Let $n=[E:F]$. Frobenius is an isomorphism from $E$ onto $E^p$ carrying $F$ onto $F^p$, so $[E^p:F^p]=n$. If $F^p=F$, then $F\subseteq E^p\subseteq E$ and both $[E^p:F]$ and $[E:F]$ equal $n$, forcing $E^p=E$.
Conversely, if $E^p=E$, the equality above reads $[E:F^p]=n$. In the tower $F^p\subseteq F\subseteq E$, it follows that $n=n[F:F^p]$, so $[F:F^p]=1$ and $F^p=F$. Finiteness of $n$ is essential for this degree argument.
